
\chapter*{Abstract}

Small language models can run on smartphones, but remembering long interactions strains their
limited context, battery, memory and thermal headroom. The on-device studies reviewed here do not compare memory
representations for a conversation history by both answer quality and energy on one handset.
This thesis examines how memory strategies affect answer quality, what resource trade-offs
they create on a phone, and where knowledge-graph memory succeeds or fails.

A shared pipeline compares truncated full context, vector retrieval and graphs extracted by
small language models. The main study runs on a Samsung Galaxy S23 Ultra and measures accuracy,
energy, latency, memory use and temperature, reporting accuracy and energy separately and as
correct answers per kilojoule. A workstation runs the same modules with more questions and
Gemma and Qwen readers on LoCoMo and LongMemEval, and tests retrieval depth, graph repairs and
source-text restoration.

On the phone, vector retrieval answered more questions correctly than truncated full context
and used less energy per correct answer. On the workstation, full context with a long window had higher
accuracy point estimates than vector retrieval on LoCoMo and lower ones on LongMemEval-S, while
using many times more prompt tokens. In a single construction pilot on the phone, building
a graph cost substantially more energy than building the vector index. The accuracy of
graph memory was measured on the workstation. There, graph memory was limited by extraction
failures and by how its evidence reached the reader, and restoring the source text behind
retrieved facts improved graph answers for most readers without reaching vector retrieval.
Selective Provenance Expansion, developed to choose between compact facts and source text under
a token budget, did not reliably recover that benefit.

The results suggest evaluating memory by both the evidence it delivers and the resources needed
to build and use it. They are limited by one handset, a small phone sample and automatic
judging, and cover question answering rather than complete mobile-agent workflows.

\medskip
\noindent\textbf{Keywords:} small language models; long-term memory; retrieval-augmented
generation; knowledge graphs; on-device inference; energy efficiency.
